\documentclass[a4paper,12pt]{article}
\usepackage[utf8]{inputenc}
\usepackage[polish]{babel}
\usepackage[T1]{fontenc}
\usepackage{geometry}
\usepackage{graphicx}
\usepackage{enumitem}
\usepackage{hyperref}
\usepackage{longtable}
\geometry{margin=2.5cm}

\title{Sprawozdanie z realizacji projektu - Język Java}
\author{Eryk Witkowski, Bartosz Badura}
\date{\today}

\begin{document}

\maketitle

\section*{1. Tytuł projektu}
\textbf{Analiza ruchu sieciowego [NetDolphin?]}\\
Aplikacja pozwalająca na monitorowanie i wykrywanie podstawowych anomalii w ruchu sieciowym

\section*{2. Założone vs. zrealizowane funkcjonalności}

\begin{longtable}{|p{6cm}|p{4cm}|p{4cm}|}
\hline
\textbf{Funkcjonalność} & \textbf{Założona} & \textbf{Zrealizowana} \\
\hline
Przechwytywanie i analiza pakietów & Tak & Tak \\
\hline
Wykrywanie podejrzanych pakietów & Tak & Tak \\
\hline
generowanie raportów bezpieczeństwa & Tak & Tak \\
\hline
\end{longtable}

\textbf{Zrealizowano: 3 z 3 zaplanowanych funkcjonalności ( 100\%)}

\section*{3. Realizacja wymagań}

\subsection*{Funkcjonalne vs. niefunkcjonalne}

\begin{longtable}{|p{6cm}|p{4cm}|p{4cm}|}
\hline
\textbf{Typ Wymagań} & \textbf{Założone} & \textbf{Zrealizowane} \\
\hline
Funkcjonalne & 8 & 7 \\
\hline
Niefunkcjonalne & 4 & 3 \\
\hline
\end{longtable}

\textbf{Procent realizacji:} Funkcjonalne: 87.5\%, Niefunkcjonalne: 75\%

\section*{4. Rzeczywisty nakład pracy}

\begin{itemize}
    \item Planowany czas pracy: 120-160 godzin
    \item Rzeczywisty czas pracy: 100 godzin
\end{itemize}

\section*{5. Napotkane problemy i rozwiązania}

\begin{itemize}
    \item \textbf{Problemy z przechowywaniem odebranych pakietów:} przechowywanie wszystkich pakietów w prostej strukturze jak ArrayList w programie powodowało zacinanie się programu, a przy większej liczbie pakietów (np. 5000) - całkowite zatrzymanie systemu\\
    \textit{Rozwiązanie:} Zmiana podejścia na dodawanie pakietów do cache'a (pliku .pcap) - pozwala na bezproblemowe przechowanie do miliona pakietów, z zastrzeżeniem co do znacząco rosnącego rozmiaru pliku

    \item \textbf{Trudność w testowaniu rozwiązania:} przez charakter rozwiązania napisanie testów jednostkowych do większości funkcjonalności było bardzo skomplikowane\\
    \textit{Rozwiązanie:} napisanie testów jedynie do niezależnych od wyłapywanych pakietów metod i klas pomocniczych
\end{itemize}

\section*{6. Możliwości Rozwoju}

\begin{itemize}
    \item Dodanie oznaczenia pakietów podejrzanych o bycie anomalią na czerwono w tabeli  
    \item Dodanie większej liczby obsługiwanych protokołów
    \item Dodanie analizy pakietów bazującej na machine-learningu
    \item Możliwość otworzenia podglądu zapisanej listy pakietów
    \item Możliwość zatrzymania wyłapywania pakietów z poziomu UI
    \item Polepszenie filtrowania pakietów - dodanie interpretera, podpowiadanie opcji/składni itp. (w tym momencie składnia jest restrykcyjna i wymaga oddzielenia precyzyjnych nazw (srcip, srcmac itd..) dwoma znakami \&)
\end{itemize}

\section*{7. Zrzuty Ekranu}
\begin{figure}[h!]
    \centering
    \includegraphics[width=0.9\textwidth]{screenshot.png}
    \caption{Wybór interfejsu sieciowego}
\end{figure}

\begin{figure}[h!]
    \centering
    \includegraphics[width=0.9\textwidth]{screenshot1.png}
    \caption{Podgląd pakietów}
\end{figure}
\begin{figure}[h!]
    \centering
    \includegraphics[width=0.9\textwidth]{screenshot2.png}
    \caption{Raport HTML}
\end{figure}

\end{document}
